\section{The fixed-order comparison on compact time intervals}
\label{sec:comparison}
\subsection{Integrate the polynomial remainder before bounding the window}
For \(m\ge0\) define
\begin{subequations}\label{eq:gaussian-budgets}
\begin{align}
 M_m&=\sqrt2\,\e^{R+Y^2/(4\tau)+2B}
   \sum_{k=0}^m\binom{m}{k}W^{m-k}(8B)^{k/2}
                        \frac{\Gamma((k+1)/2)}{\sqrt\pi},\\
 M_m^t&=\frac{M_m}{2\tau}+\frac{M_{m+2}}{4\tau^2}.
\end{align}
\end{subequations}
Dependence on \(\tau,B,r,b\) is suppressed in these symbols.
\begin{lemma}[Complex Gaussian moments]\label{lem:moments}
For \(z\in\Omega_X(R,Y)\), \(\tau\le T\le B\), and \(q=u-X\),
\begin{subequations}\label{eq:moment-bounds}
\begin{align}
 \int_\R\abs{k_T(z-u)}\e^{\abs q}\abs q^m\dd u&\le M_m,\\
 \int_\R\abs{\partial_Tk_T(z-u)}\e^{\abs q}\abs q^m\dd u&\le M_m^t.
\end{align}
\end{subequations}
\end{lemma}
\begin{proof}
Put \(v=u-\Rea z\).
Then \(\abs q\le R+\abs v\) and
\(\max(\abs q,\abs{z-u})\le W+\abs v\).
The inequality \(\abs v\le v^2/(8T)+2T\) leaves the Gaussian
\(\e^{-v^2/(8T)}\).
Its \(k\)-th absolute moment, including the original kernel
normalization, is
\(\sqrt2(8T)^{k/2}\Gamma((k+1)/2)/\sqrt\pi\).
The binomial theorem proves the first inequality.
For the second differentiate the actual kernel:
\begin{equation}
 \partial_Tk_T(z-u)
 =\left(\frac{(z-u)^2}{4T^2}-\frac1{2T}\right)k_T(z-u).
 \label{eq:time-kernel}
\end{equation}
The same moment calculation proves \eqref{eq:gaussian-budgets},
including both inverse-time factors.
\end{proof}

At the real center \(\abs{b_{\epsilon n}}=g n^{-1/2}\).
Sectorial calibration gives \(\abs{\Rea\lambda_{\epsilon n}}\le1\)
after a fixed onset. Therefore
\begin{equation}
 \sum_{\epsilon=\pm1}\sum_{n=1}^N
       \abs{b_{\epsilon n}\e^{\lambda_{\epsilon n}q}}
 \le C_{\mathrm{pop}}gX^{1/4}\e^{\abs q},
 \qquad C_{\mathrm{pop}}=4(4\pi)^{-1/4}.
 \label{eq:population}
\end{equation}
Here \(\sum_{n\le N}n^{-1/2}\le2\sqrt N\).
Combining \cref{lem:weighted,lem:moments} with
\eqref{eq:population} gives local field and direct-time errors
\begin{equation}
 C_{\mathrm{pop}}\e\,gX^{-(p+3/4)}\mathsf B_{p+1}[M],
 \qquad
 C_{\mathrm{pop}}\e\,gX^{-(p+3/4)}\mathsf B_{p+1}[M^t].
 \label{eq:local-payments}
\end{equation}
The time bound uses moments through degree \(2p+4\).
No factor \((\log X)^{2p+2}\) appears: the powers of \(q\) were
integrated, not replaced by \(V_B\).

\subsection{Why the completion order must track the curvature order}
\label{subsec:matching}
On the window, sectorial gamma estimates give
\(G(X+q)\le C_\Gamma g\e^{\abs q}\).
If \(X-V_B\ge X/2\), \eqref{eq:rs-remainder} becomes
\begin{subequations}\label{eq:CK}
\begin{align}
 \abs{H(0,u)-f_K(u)}
 &\le C_Kg\e^{\abs q}X^{-(K/2+3/4)},\\
 C_K&=C_\Gamma\frac{2^{7/4}}7
 \Gamma((K+1)/2)1.1^{K+1}(8\pi)^{K/2+3/4}.
\end{align}
\end{subequations}
Its field and direct-time integrals cost \(C_KM_0\) and \(C_KM_0^t\)
at this power. Thus the local curvature remainder and the canonical
initial remainder have the same exponent precisely when
\begin{equation}
 K=2p.
 \label{eq:matching}
\end{equation}
This is the least integer order supporting that rate by the
displayed upper bound. It is not a lower bound on the actual
Riemann--Siegel error. Retaining only \(K=2p-1\) leaves an
\(X^{-(p+1/4)}\) allowance even though the arithmetic approximation
has improved further.

\subsection{True tails and artificial tails}
The whole-line polynomial transform in \cref{sec:hierarchy} creates
an artificial tail outside the local window. It is distinct from
the true initial tail of \eqref{eq:heat-convolution}.
Introduce
\begin{subequations}\label{eq:tail-moments}
\begin{align}
 \mathsf D_m&=\frac{\e^{Y^2/(4\tau)}}{\sqrt{4\pi\tau}}\,
      2^m(64B)^{(m+1)/2}\Gamma((m+1)/2),&
 \mathsf D_m^t&=\frac{\mathsf D_m}{2\tau}
                        +\frac{\mathsf D_{m+2}}{4\tau^2},\\
 \mathsf J_m&=\frac{\e^{Y^2/(4\tau)}}{\sqrt{4\pi\tau}}
  \sum_{k=0}^m\binom{m}{k}Y^{m-k}
                 (8B)^{(k+1)/2}\Gamma((k+1)/2),&
 \mathsf J_m^t&=\frac{\mathsf J_m}{2\tau}
                        +\frac{\mathsf J_{m+2}}{4\tau^2}.
\end{align}
\end{subequations}
When \(V_B\ge\max(2R,W+1,32B)\), for \(\abs q\ge V_B\),
\[
 -\frac{(q-\Rea d)^2}{4B}+\abs q\le-\frac{q^2}{32B},
 \qquad \abs{z-u}\le2\abs q.
\]
Splitting the last Gaussian in half and extracting its value at
\(V_B\) gives the tail moment bounds
\(\mathsf D_m\e^{-V_B^2/(64B)}\le\mathsf D_m\e^{-L^2}\)
and their direct-time versions.

The real canonical initial field has the height-sensitive estimate
\[
 \abs{H(0,X+q)}
 \le C X^{P_0}\e^{-\pi X/8+\pi\abs q/8},
 \qquad \abs q\le X/2,\qquad P_0=4.
\]
It follows from the gamma exponential and a polynomial vertical-strip
bound for zeta; the deliberately coarse exponent \(4\) is sufficient.
Only beyond \(\abs q>X/2\) do we use the global bound
\(\abs{H(0,u)}\le M_H=\int_0^\infty\abs{\Phi(v)}\dd v\).
For \(X\ge4R\), that region has
\(\abs{u-\Rea z}\ge X/4\). Splitting its Gaussian gives the
remote factor \(\e^{-X^2/(128B)}\), and division by \(g\)
costs \(\e^{\pi X/8}X^{-7/4}/g_-\).

Absorbing only fixed source constants into \(C_{\mathrm{tail}}\), set
\begin{subequations}\label{eq:tail-constants}
\begin{align}
 A_{\mathrm{mod}}&=C_{\mathrm{tail}}
                 \max(\mathsf D_0,\mathsf D_0^t),\\
 A_{\mathrm{art},p}&=C_{\mathrm{pop}}
       \max\{\mathsf H_p[\mathsf D],\mathsf H_p[\mathsf D^t]\},\\
 A_{\mathrm{far}}&=(M_H/g_-)\max(\mathsf J_0,\mathsf J_0^t).
\end{align}
\end{subequations}
The following bounds hold for both the field and the directly
computed time derivative, after division by \(g\):
\begin{center}
\begin{tabular}{@{}ll@{}}
\toprule
Contribution & Upper bound\\
\midrule
True data, \(V_B<\abs q\le X/2\) &
\(A_{\mathrm{mod}}X^{P_0-7/4}\e^{-L^2}\)\\
Artificial polynomial data, \(\abs q>V_B\) &
\(A_{\mathrm{art},p}X^{1/4}\e^{-L^2}\)\\
True data, \(\abs q>X/2\) &
\(A_{\mathrm{far}}X^{-7/4}\e^{\pi X/8-X^2/(128B)}\)\\
\bottomrule
\end{tabular}
\end{center}
No gamma function or Riemann--Siegel function is extended onto
the artificial tails.

\subsection{The theorem and its derivation}
\begin{theorem}[Fixed order on a fixed compact time family]
\label{thm:fixed}
Fix \(p\ge1\), \(r,b>0\), and \(0<\tau\le B<\infty\).
For \(K=2p\), the model \eqref{eq:model} satisfies, for all
sufficiently large \(X\),
\begin{equation}
 \normg{h-F_{X,B}^{[p,2p]}}
 \le C_{p,r,b}(\tau,B)X^{-(p+3/4)}.
 \label{eq:fixed-result}
\end{equation}
A sufficient constant is
\begin{subequations}\label{eq:fixed-constant}
\begin{align}
 A_{f,p}&=C_{2p}M_0+C_{\mathrm{pop}}\e\,\mathsf B_{p+1}[M]+3,\\
 A_{t,p}&=C_{2p}M_0^t+C_{\mathrm{pop}}\e\,\mathsf B_{p+1}[M^t]+3,\\
 C_{p,r,b}(\tau,B)&=C_\chi(48A_{f,p}+A_{t,p}).
\end{align}
\end{subequations}
The sufficient onset is specified in \cref{app:onsets}.
The frozen heat residual and its first spatial derivative vanish.
\end{theorem}
\begin{proof}
Split \eqref{eq:heat-convolution} at \(I_{X,B}\), and on that
window insert
\[
 H(0,u)=S_N(u)+Q_{X,2p}(u)+[H(0,u)-f_{2p}(u)].
\]
The \(Q\)-integral cancels exactly against the retained completion.
For each mode subtract the whole-line polynomial transform.
The remaining local difference is bounded by
\eqref{eq:local-payments}; the initial mismatch is bounded by
\eqref{eq:CK}; the three tails are bounded above.
The onset inequalities in \cref{app:onsets} put each tail below
\(X^{-(p+3/4)}\). This proves the field bound \(gA_{f,p}X^{-(p+3/4)}\)
on the enlarged chart and, independently, the time bound
\(gA_{t,p}X^{-(p+3/4)}\).

Both fields are holomorphic in \(z\). Radius-\(1/2\) Cauchy disks
fit inside the enlarged chart, and give factors \(j!2^j\le48\)
for \(0\le j\le3\). Convert from \(g\) to \(\abs{G(z)}\)
using \(C_\chi\). This proves \eqref{eq:fixed-constant}.
The residual assertion follows from \cref{prop:residual}.
\end{proof}
The theorem includes a one-point time interval \(\tau=B\), since
the underlying smooth functions are defined for \(T>0\).
Its time derivative is not a derivative of an error allowance.
The explicit constants deteriorate through
\(\e^{Y^2/(4\tau)+2B}\), powers of \(B\), and
\(\tau^{-1},\tau^{-2}\). No uniform limit at \(\tau=0\)
or \(B=\infty\) is asserted.

\subsection{Low-order formulas and the earlier first correction}
\label{subsec:low-orders}
Write \(\gamma_\epsilon=\ell_\epsilon''(X)\) and
\(\eta_\epsilon=\ell_\epsilon'''(X)\). Then
\[
 U_{\epsilon,1}=\frac{\gamma_\epsilon}{2}\calQ_T^2U_{\epsilon,0},
 \quad
 U_{\epsilon,2}=\left(\frac{\eta_\epsilon}{6}\calQ_T^3+
             \frac{\gamma_\epsilon^2}{8}\calQ_T^4\right)U_{\epsilon,0}.
\]
For one mode the first corrected expression is
\begin{equation}
 b\e^{\lambda d+T\lambda^2}
       \left[1+\frac{\gamma}{2}\{(d+2T\lambda)^2+2T\}\right].
 \label{eq:first-frozen}
\end{equation}
At \(d=0\), its bracket is
\(1+\gamma(T+2T^2\lambda^2)\).
This agrees with the arithmetic value of the earlier
target-dependent correction \cite{ans10},
\[
 h_{\epsilon n}(z)\e^{T\beta_{\epsilon n}(z)^2}
    [1+\gamma_\epsilon(z)(T+2T^2\beta_{\epsilon n}(z)^2)],
 \quad
 \beta_{\epsilon n}(z)=\partial_z\log h_{\epsilon n}(z),
\]
evaluated at \(z=X\). In that formula, however, the gamma data
vary with the target. Its spatial derivatives differ from those
of \eqref{eq:first-frozen}, and its residual is
\(O(\abs G X^{-7/4})\), not identically zero.
Its \(K=1\) completion gives a total \(X^{-5/4}\) error allowance.
For the frozen model, the local first-level allowance is
\(X^{-7/4}\); retaining \(K=2\) supports it. At the second level
the mixed curvature-square term and \(C_0,\ldots,C_4\) must all
remain to support \(X^{-11/4}\).

Thus higher approximation means retaining all partitions of the
chosen fluctuation weight and matching the completion order.
Adding only one higher derivative of gamma does not implement
the next level.
